\hypertarget{the-audit-account}{%
\chapter{19. The audit account}\label{the-audit-account}}

A reference monitor is asked for three things. It must mediate every
access, it must resist tampering, and it must be small enough to verify
and honest enough to account for what it did. The volume so far has been
the first two made mechanism: every crossing driven through one gateway,
the enforcer always outside the thing it confines, privilege held in the
narrowest place and never in the daemon a workload can reach. What
remains is the third, and it is two accounts rather than one. Before a
kennel runs, a report says what its policy exposes. While it runs, a log
records what it did. Together they are how the monitor answers for what
it mediates.

\hypertarget{the-risk-report}{%
\section{19.1 The risk report}\label{the-risk-report}}

\texttt{kennel\ policy\ risks} evaluates a fully folded policy against
the threat catalogue and reports what the grants expose and what they
mitigate. Each finding names a threat by its catalogue identity, carries
the catalogue's one-line residual for it, names the granting site that
carries it, and quotes the reason the operator wrote there. So an
exposure is never a bare label: it is a threat, the exact grant that
opens it, the residual that remains after the framework's mitigation,
and the operator's own stated cause for accepting it, read off the
policy as a single legible line.

\begin{kennelcode}
$ kennel policy risks
exposes
  T1.6  host-network reconnaissance, the residual after host mode
        at   [net] mode = "host"
        why  "the build needs the corporate proxy on a host-only VLAN"
mitigates
  T3.8  arbitrary egress, closed by the proxy default-deny
\end{kennelcode}

An exposure reads as the threat, the residual the catalogue records for
it, the exact grant that opens it, and the reason the author wrote at
that grant, so the report answers not only what a policy risks but why
its author accepted the risk.

A tag reaches the report two ways, and the report keeps them apart. Most
are authored, written on a grant by whoever wrote the template or the
leaf. Some are derived: the compiler infers them from the shape of a
grant, so a policy that opens host network mode has the
host-reconnaissance residual reinstated against it whether or not the
author tagged it (T1.6), a device passthrough carries the
control-deactivation exposure (T2.1), a headless signing allowance
reinstates the lateral-reach residual (T1.6). The derived set is what
makes the report match what is enforced rather than only what someone
remembered to write, and marking each finding authored or derived is
what keeps an inference the compiler drew from being mistaken for a
claim the author made.

The honesty of the report is structural rather than asserted. The threat
tags live only in the source policy and never survive into the settled
artefact, so the report reads the same folded source the compiler
enforced from, which is the honest input and not a separately maintained
story about it. A tag that names no catalogued threat is surfaced as a
likely typo rather than silently dropped or counted as a real risk, and
a policy authored against an older catalogue version earns a drift note.
There is no invented risk and no theatre: the report says what the
grants carry, what the compiler can infer, and what it could not
resolve, and nothing beyond that. This is the warning surface the rest
of the volume kept pointing at (\texttt{footgun-warn-dont-forbid}). Host
network mode's reinstated residual, the delegated-spawning exposure, the
substrate trust of a pulled image, the standing-service consume: each
was called a loud grant surfaced in the risk report, and this is the
report, the place a deliberate exposure is shown to the operator as the
reasoned choice it was rather than forbidden outright or buried.

\hypertarget{the-audit-log}{%
\section{19.2 The audit log}\label{the-audit-log}}

While a kennel runs, the monitor records what it mediated in a
structured per-kennel log. The record is tunable by resource class, the
network, the filesystem, execution, the local-socket surface, and the
bus each carrying their own setting, and by level within a class, from
off through denies-only and a per-class summary to the full record of
every mediated access. An operator logs every filesystem access and only
denials on the network, or the reverse, or everything, by the
deployment's need rather than a single global switch. The log is written
to one of a few sinks, per-class files of newline-delimited JSON by
default, or the system journal, or syslog, or the daemon's own output
for a container that collects it from there.

\begin{kennelcode}
{"schema_version":1,"ts":"2026-06-29T12:37:33.829451Z","kennel":"build-agent",
 "kennel_uuid":"7f3c...","event":"net.connect","resource":"net","outcome":"deny",
 "source":"kenneld","host":"remco-vm","pid":58471,"dest":"exfil.example:443"}
\end{kennelcode}

Every record carries the same envelope, the schema version, a timestamp,
the kennel name and its run identity, the event and resource class, the
outcome, and the daemon as source, with the resource-specific fields
beside it. A denial names what was refused and a mediation what was
allowed, so the log is the account of what the monitor did, not what the
policy said it might.

One definition governs all of this, shared between the compiler that
reads a policy's audit section and the daemon that reads a standalone
defaults file, and it lives in the runtime crate rather than the
compiler. That placement is the same discipline that kept the resolver
out of the daemon, applied once more: the daemon must resolve audit
configuration at spawn time, so the schema it needs sits where it can
reach it without linking the compiler the daemon is built to exclude.
The format carries a schema version and is stable across releases, so a
tool written to parse the log can rely on its structure holding.

\hypertarget{what-the-account-is-for-and-where-it-stops}{%
\section{19.3 What the account is for, and where it
stops}\label{what-the-account-is-for-and-where-it-stops}}

The two accounts carry the uses that justify keeping them. A log of
every resource access answers a forensic question after the fact, what a
workload read between one time and another. It is evidence for a
compliance claim, because a kennel that did not reach a protected
resource did not reach it by the kernel's refusal and the log holds the
refusals. It feeds policy refinement, a denied access the workload
legitimately needed becoming the next grant the operator adds with a
reason. And it is the input to anomaly detection, a monitoring layer
watching for new destinations or unusual access or repeated denials.
Every one of those is someone else's to build on top: the framework
provides the record, never the watcher over it.

That boundary is the honest shape of the whole account. The log is
write-once and append-only from the framework's side, and stronger
tamper-evidence, signed segments or an external transparency log, is a
downstream concern the format leaves room for without claiming to
provide. The compliance story is bounded the same way: the framework
claims no regulatory certification and provides the infrastructure that
makes a certification-relevant claim demonstrable, the evidence rather
than the verdict. The risk report says what a policy exposes and the log
says what a kennel did; whether either is acceptable is a judgement with
standing the framework does not assume for itself. It accounts; it does
not absolve.

This is where the volume ends, because it is where the reference
monitor's third demand is met. The account does not strengthen the
confinement, which holds whether or not a log is ever read and whether
or not a risk report is ever run. It makes the confinement answerable.
What a policy exposed can be said before it runs, what a kennel did can
be read after, and both can be checked against the design the first
volume set out, in a form a person or a tool can verify rather than take
on faith. A monitor that mediates completely and resists tampering but
cannot be inspected still asks to be trusted. One that also accounts for
what it mediated, in the open and in its own grant's words, asks for
less, and that is the whole of what this design has tried to do: hold
the line, and show its working.
